\documentclass[11pt]{article}
\usepackage[margin=2.1cm]{geometry}
\usepackage{booktabs}
\usepackage{amsmath}
\usepackage{enumitem}
\usepackage[T1]{fontenc}
\usepackage{lmodern}

\setlength{\parindent}{0pt}
\setlength{\parskip}{5pt}
\newcommand{\code}[1]{\texttt{#1}}
\newcommand{\answerbox}[1]{\vspace{2pt}\fbox{\parbox[c][#1][c]{\dimexpr\linewidth-2\fboxsep-2\fboxrule}{\ }}\vspace{4pt}}

\begin{document}

\begin{center}
{\Large\bfseries EE\,362 --- Pre-Lab 11}\\[2pt]
{\large Instruction Decoder II}\\[6pt]
\begin{tabular}{@{}p{6.5cm}p{4cm}p{4cm}@{}}
Name: \hrulefill & Section: Mon / Tue & Date: \hrulefill
\end{tabular}
\end{center}

\vspace{2pt}\hrule\vspace{4pt}

{\small\textbf{Due at the start of the lab session.} Five points, about twenty
minutes, on paper.}

\vspace{4pt}\hrule

\section*{Question 1 --- The four new rows}

Fill in only the signals that differ from \code{CTRL\_DEFAULT}.

\begin{center}
\small
\begin{tabular}{lllllll}
\toprule
instr & \code{ALU\_Asel} & \code{IR\_Xlsft} & \code{DM\_en} & \code{DM\_rw} & \code{RF\_strb} & \code{RF\_oe} \\
\midrule
\code{LW}   & \rule{1.6cm}{0.4pt} & \rule{1.4cm}{0.4pt} & \rule{1.4cm}{0.4pt} & \rule{1.2cm}{0.4pt} & \rule{1.6cm}{0.4pt} & \rule{1.4cm}{0.4pt} \\[6pt]
\code{SW}   & \rule{1.6cm}{0.4pt} & \rule{1.4cm}{0.4pt} & \rule{1.4cm}{0.4pt} & \rule{1.2cm}{0.4pt} & \rule{1.6cm}{0.4pt} & \rule{1.4cm}{0.4pt} \\[6pt]
\code{BLEZ} & \rule{1.6cm}{0.4pt} & \rule{1.4cm}{0.4pt} & \rule{1.4cm}{0.4pt} & \rule{1.2cm}{0.4pt} & \rule{1.6cm}{0.4pt} & \rule{1.4cm}{0.4pt} \\[6pt]
\code{J}    & \rule{1.6cm}{0.4pt} & \rule{1.4cm}{0.4pt} & \rule{1.4cm}{0.4pt} & \rule{1.2cm}{0.4pt} & \rule{1.6cm}{0.4pt} & \rule{1.4cm}{0.4pt} \\
\bottomrule
\end{tabular}
\end{center}

\textbf{(b)} \code{SW} is the only instruction with \code{RF\_oe = DRIVE}.
Explain what would physically happen on the memory data bus if it were left at
\code{HiZ}, and what would end up stored.

\answerbox{2.2cm}

\textbf{(c)} \code{SW}, \code{BLEZ} and \code{J} all have
\code{RF\_strb = HOLD}. State the general rule in one sentence: which
instructions strobe the register file, and which do not?

\answerbox{1.6cm}

\section*{Question 2 --- Reading the completed table}

\textbf{(a)} \code{ALU\_Asel = IMMED} for \code{ADDI}, \code{LUI}, \code{LW} and
\code{SW}. Those four look unrelated. What do they actually have in common, and
what does that tell you about the hardware?

\answerbox{2.4cm}

\textbf{(b)} \code{BLEZ} sets \code{ALU\_func = ADD} even though a branch does
not compute a value anybody keeps. Why does the ALU run at all?

\answerbox{2.0cm}

\section*{Question 3 --- Don't cares, and testing}

A classmate implements \code{IR\_Xsel} as ``SIGN for everything except
\code{LUI}''. You implemented it as ``SIGN for \code{ADDI}, \code{LW},
\code{SW}, \code{BLEZ}''.

\textbf{(a)} On which instructions do your two designs give different answers?

\answerbox{1.4cm}

\textbf{(b)} Is either of you wrong? Justify your answer in terms of which
instructions actually read the extender output.

\answerbox{2.2cm}

\textbf{(c)} Your testbench checks \code{IR\_Xsel} for all eight instructions,
including \code{ADD}. What is wrong with that, and what should it check instead?

\answerbox{2.2cm}

\vfill
{\small\emph{Question 3 is the one that matters.} It is the same question as
Section 5 of your lab report.}

\end{document}
